\documentclass[11pt,a4paper]{article}
\usepackage{amsmath,amssymb,amsthm}
\usepackage[margin=2.5cm]{geometry}
\usepackage{hyperref}

\newtheorem{definition}{Definition}[section]
\newtheorem{theorem}[definition]{Theorem}
\newtheorem{proposition}[definition]{Proposition}
\newtheorem{corollary}[definition]{Corollary}
\newtheorem{remark}[definition]{Remark}
\newtheorem{conjecture}[definition]{Conjecture}

\title{Hyperseed Formalization:\\
  Announcing the Beneficial AGI Summit}
\author{ProtomegaTron (automated formalization)}
\date{2026-09-19 (deepened v2)}

\begin{document}
\maketitle

\begin{abstract}
We formalize Ben Goertzel's ``Announcing the Beneficial AGI Summit''
(2023-10-25) within the Hyperseed ontology, incorporating d-calculus
extensions from Hyperseed v2. The summit is modeled as a fiber
cross-pollination event, the unconference format as emergent fiber
structure, beneficial AGI as fiber-ethics coupling, and key concepts
(trust, governance, diversity) are given d-calculus formulations as
connection coefficients, holonomy, and portfolio curvature.
\end{abstract}

\tableofcontents

\section{Summit as fiber cross-pollination}

\begin{definition}[Research fiber bundle]
Each research group $g$ develops a fiber bundle
$\pi_g: E_g \to B_g$ where $B_g$ is the group's research context
(problems, data, methods) and $E_g$ is the group's theoretical fiber
(concepts, results, frameworks).
\end{definition}

\begin{proposition}[Cross-pollination event --- claims 1, 5, 8]
\label{prop:cross}
The summit creates a temporary shared base space
$B_{\text{summit}} \supset \bigcup_g B_g$ where fibers from different
groups overlap. At overlap points $b \in B_g \cap B_h$:
\begin{enumerate}
  \item \textbf{Transfer:} Fiber elements from $E_g(b)$ can be mapped
    to $E_h(b)$ via a cross-pollination morphism $\chi_{gh}$.
  \item \textbf{Compatibility:} Sections $\sigma_g$ and $\sigma_h$
    can be compared at $b$: do they agree? contradict? complement?
  \item \textbf{Emergence:} New fiber $E_{\text{new}}(b)$ emerges at
    contact points that didn't exist in either $E_g$ or $E_h$.
\end{enumerate}
The emergent fiber is the cognitive synergy of the research community.
\end{proposition}

\section{Unconference as emergent fiber}

\begin{proposition}[Emergent vs imposed structure --- claim 4, 9]
\label{prop:emergent}
Traditional conferences impose fiber structure (fixed schedule, invited
talks). Unconferences let fiber structure emerge:
\begin{itemize}
  \item \textbf{Imposed:} $E_{\text{conf}} = E_{\text{organizer}}$
    (fiber determined by organizers).
  \item \textbf{Emergent:} $E_{\text{unconf}} = \text{colim}_g(E_g)$
    (fiber is the colimit of participant fibers --- the minimal
    structure compatible with all participant contributions).
\end{itemize}
The emergent fiber is better adapted because it's shaped by actual
participant fiber structure rather than organizer assumptions.
\end{proposition}

\section{Beneficial AGI as fiber-ethics coupling}

\begin{proposition}[Beneficence coupling --- claims 2, 7, 10]
\label{prop:beneficial}
``Beneficial AGI'' imposes a coupling constraint on fiber evolution:
\[
  \frac{dF}{dt} > 0 \quad \text{AND} \quad
  \frac{d\,\mathcal{W}}{dt} > 0
\]
where $F$ is capability fiber and $\mathcal{W}$ is flourishing
(well-being) fiber. Capability growth without flourishing growth is
dangerous; flourishing growth without capability is ineffective.
The coupling ensures they co-evolve.

Corporate AGI optimizes $dF/dt$ without the $d\mathcal{W}/dt$ constraint.
The BGI movement provides competitive pressure to enforce the coupling.
\end{proposition}

\section{d-calculus formulation (Hyperseed v2)}

\begin{proposition}[Cross-pollination curvature --- claim 11]
\label{prop:curvature}
The curvature $F_\nabla^{gh}$ at fiber contact points between groups
$g$ and $h$ measures cross-pollination difficulty:
\[
  \|F_\nabla^{gh}\| \ll 1 \implies \text{easy transfer (compatible frameworks)}
\]
\[
  \|F_\nabla^{gh}\| \gg 1 \implies \text{difficult transfer (incompatible,
    but potentially more valuable)}
\]
High-curvature cross-pollination is harder but produces more novel
emergent fiber --- the cognitive synergy benefit is proportional to
curvature overcome.
\end{proposition}

\begin{proposition}[Emergent fiber gradient --- claim 12]
\label{prop:gradient}
The emergent fiber at the summit has a gradient $\nabla F_{\text{emergent}}$
that points in directions no individual group was pursuing:
\[
  \nabla F_{\text{emergent}} \notin \text{span}\{\nabla F_g : g \in \text{groups}\}
\]
This gradient represents genuinely new research directions that arise
only from collective interaction.
\end{proposition}

\begin{definition}[Trust as connection coefficient --- claim 13]
\emph{Trust} between researchers is a connection coefficient
$\Gamma_{\text{trust}}^{ij}$ in the social fiber bundle. It enables
parallel transport of sensitive information (safety concerns, alignment
strategies, vulnerability disclosures):
\[
  \Gamma_{\text{trust}}^{ij} > \theta \implies
  \text{sensitive fiber can transport from } i \text{ to } j
\]
Building trust at the summit = increasing $\Gamma_{\text{trust}}^{ij}$
for researcher pairs $(i,j)$.
\end{definition}

\begin{proposition}[Governance holonomy --- claims 3, 14]
\label{prop:holonomy}
Decentralized governance cycles at the summit
\[
  \gamma: \text{plan} \to \text{meet} \to \text{discuss}
    \to \text{decide} \to \text{implement} \to \text{evaluate}
\]
produce non-trivial holonomy in the governance fiber bundle:
\[
  \operatorname{Hol}_\gamma(\nabla) \neq \mathrm{id}
  \implies \text{governance evolves with each cycle}
\]
This is the mark of a living governance structure, in contrast to
static regulatory frameworks that produce trivial holonomy (cycling
without evolution).
\end{proposition}

\begin{proposition}[Portfolio curvature --- claim 15]
\label{prop:portfolio}
The total curvature across the portfolio of AGI approaches measures
diversification benefit:
\[
  \kappa_{\text{portfolio}} = \sum_{g < h} \|F_\nabla^{gh}\|
\]
\begin{itemize}
  \item Low $\kappa_{\text{portfolio}}$: approaches are too similar
    (concentrated risk, low curvature between them).
  \item High $\kappa_{\text{portfolio}}$: approaches are genuinely
    diverse (distributed risk, high curvature between them).
\end{itemize}
The summit maximizes $\kappa_{\text{portfolio}}$ by bringing together
maximally diverse approaches --- this is portfolio diversification
formalized as curvature maximization.
\end{proposition}



%======================================================================
\section{OCO/2 crosswalk (oco/2:2.0.0-alpha.1)}
\emph{Added 2026-10-03. Interpretive mapping onto the Omega Core Ontology
record registry; OCO/2 is an unfrozen alpha candidate.}

\begin{proposition}[Diversity and emergence --- claims 16--19]
Results from Plans with distinct origins add under origin-keyed counting. An agenda whose entries are appended by participants $p_1,\dots,p_n$ has $n$ issuers, so no one participant determines its structure.
\end{proposition}

\end{document}
